%% =====================================================================
%%  7 / 7   Abstract  摘要 + 全文统稿                           成员 7
%%  五句话结构：背景 / 空白 / 做法 / 结果 / 结论
%% =====================================================================

Medical artificial intelligence is moving from systems that recommend to systems that act. Large language model agents now plan, call tools and complete multi-step clinical tasks, yet they are evaluated almost exclusively on accuracy: the strongest public benchmark reports a success rate below 70\%. This paper asks what changes when the unit of analysis moves from the advice to the action. We conduct a purposive review of seventeen studies spanning engineering and human-factors work, and identify four themes --- deskilling, the supervision illusion, accountability diffusion, and trust under uncertainty --- which we show to be symptoms of a single missing object: an explicit authorization boundary. Existing frameworks describe how much autonomy a system has; none specifies what a delegation must contain. We propose \emph{governed delegation}, which requires four things of any agent permitted to act clinically: a stated scope of authority, revocability while the action is in flight, auditability afterwards, and retention of unaided human skill. We then specify a $2\times2$ vignette experiment that tests the advisor-to-agent shift directly. The human--AI relationship in medicine is less a question of trust than a question of who may act, and on whose authority.
